\chapter{Protocol design} \label{ch:protocol-design}
In this chapter the protocol design is presented.
The protocol leverages smart contracts and asymmetric encryption to ensure confidentiality and data integrity.
The protocol adopts a multi-layered encryption approach, performed off-chain, and uses the smart contract as a communication channel for the crowdsensing campaign.
The key exchange is made to and from the campaign initiator (referred from now as \"\texttt{B}\").
\texttt{B} uses their private/public key pair for enhancing confidentiality, in this way, only \texttt{B} can access the other actors' public key, namely senders (\texttt{S}) and verifiers (\texttt{V}).
This approach shares similarities with prior work on asymmetric key exchange mediated by smart contracts~\cite{muth2020smartdhx}, although it differs substantially in both application domain and protocol design.

\section{Actors}\label{sec:actors}
The protocol involves several actors that are described in the following subsections.
\subsection{Campaign initiator}\label{subsec:campaign-initiator}
The campaign initiator (\texttt{B}) is the actor interested for collecting data for some reason.
They define the geographic area where the smart contract works and set the smart contract hyperparameters, all configurable parameters are explained in Appendix~\ref{ch:appendix-a-parameters}.
In addition they deposit the reward funds into the smart contract to ensure that funds are available and senders and verifiers can be compensated.
The details of rewarding model are described in the~\ref{subsec:proto-rewarding-model} subsection.

\subsection{Sender}\label{subsec:sender}
\subsection{Verifier}\label{subsec:verifiers}

\section{Protocol workflow}\label{sec:proto-workflow}
\subsection{Initialization}\label{subsec:proto-init}
\subsection{Enrollment}\label{subsec:proto-enrollment}
\subsection{Data submission}\label{subsec:proto-datasubmit}
\subsection{Verification}\label{subsec:proto-verification}
\subsection{Data publishing}\label{subsec:proto-data-publish}
\subsection{Rewarding model}\label{subsec:proto-rewarding-model}
\subsection{Bond model}\label{subsec:proto-bond-model}
\section{Security analysis and issues}\label{sec:security-analysis-issues}